% SPDX-License-Identifier: CC-BY-4.0
\documentclass[11pt,letterpaper]{article}

\usepackage[T1]{fontenc}
\usepackage{amsmath,amssymb}
\usepackage{tikz}
\usepackage{pgfplots}
\pgfplotsset{compat=1.18}
\usepackage{algorithm}
\usepackage{algpseudocode}
\usepackage{booktabs}
\usepackage{enumitem}
\usepackage[hidelinks]{hyperref}

\newcommand{\N}{\mathbb{N}}
\newcommand{\Z}{\mathbb{Z}}
\newcommand{\Q}{\mathbb{Q}}
\newcommand{\R}{\mathbb{R}}
\newcommand{\C}{\mathbb{C}}
\newcommand{\Map}[3]{#1 \colon #2 \to #3}
\newcommand{\Mod}[1]{#1\text{-}\mathrm{Mod}}
\newcommand{\GroupHom}{\mathrm{Hom}_{\mathbf{Grp}}}
\newcommand{\RingHom}{\mathrm{Hom}_{\mathbf{Ring}}}

\newtheorem{theorem}{Theorem}
\newtheorem{lemma}[theorem]{Lemma}
\newtheorem{corollary}[theorem]{Corollary}
\newtheorem{definition}[theorem]{Definition}
\newtheorem{remark}[theorem]{Remark}
\newenvironment{proof}{\noindent\textit{Proof.}}{\hfill\(\square\)}

\title{Generalized Collatz Library}
\author{}
\date{\today}

\begin{document}
\maketitle
\section*{Homogeneous Scalar Dynamics on Quadratic Lattices}

\begin{abstract}
For signed scalar coefficients of magnitude at least two and zero offset, the raw map on every specified imaginary quadratic coordinate order has a complete escape--periodicity dichotomy. Multiplicatively independent magnitudes force every nonzero seed to escape in complex norm. Dependent magnitudes give a finite coordinate box containing the entire orbit, hence eventual periodicity. The coordinate gcd is the exact positive homogeneous factor; its reconstruction retains the original division guard and every raw tick.
\end{abstract}

\section{Coordinates and content}

Fix integers \(\tau,\nu\) with \(4\nu>\tau^2\), and let
\[
\theta=\frac{\tau+i\sqrt{4\nu-\tau^2}}2,\qquad R=\Z+\Z\theta.
\]
The relation \(\theta^2=\tau\theta-\nu\) makes \(R\) a ring.
The coefficients \(x,y\) of \(z=x+y\theta\) are unique because \(\operatorname{Im}\theta>0\).
For integers \(q,d\) with \(|q|,|d|\ge2\), scalar divisibility is
\[
d\mid_R (x+y\theta)\quad\Longleftrightarrow\quad d\mid x\ \hbox{and}\ d\mid y.
\]
Indeed multiplication by the real scalar \(d\) multiplies both coordinates, and uniqueness recovers each integer quotient.
Thus the zero-offset raw map is exactly
\[
(x,y)\longmapsto
\begin{cases}
(x/d,y/d),&d\mid x,\ d\mid y,\\
(qx,qy),&\text{otherwise}.
\end{cases}
\]

For a nonzero pair \(w=(x,y)\), put \(g(w)=\gcd(|x|,|y|)>0\).
B\'ezout's identity gives
\[
d\mid x,\ d\mid y\quad\Longleftrightarrow\quad |d|\mid g(w).
\]
Also \(g(qx,qy)=|q|g(w)\), and at a legal division
\(g(x/d,y/d)=g(w)/|d|\).
Consequently the content follows the positive scalar map
\[
g\longmapsto
\begin{cases}
g/|d|,&|d|\mid g,\\
|q|g,&\text{otherwise}.
\end{cases}
\]
The pair stays nonzero because both branch multipliers are nonzero.

\begin{lemma}[Coordinate reconstruction]
\label{lem:coordinates}
If \(w_j=(x_j,y_j)\) and \(g_j=g(w_j)\), then for every \(j\ge0\),
\[
g_0|x_j|=g_j|x_0|,\qquad g_0|y_j|=g_j|y_0|.
\]
\end{lemma}

\begin{proof}
At an affine step, all three magnitudes \(g,|x|,|y|\) are multiplied by \(|q|\).
At a division step, all are divided exactly by \(|d|\).
Multiplying the proposed next identity by the nonzero integer \(|d|\) reduces it to the previous identity.
Both identities hold at time zero, so induction proves them at every raw time.
\end{proof}

\section{The classification}

\begin{theorem}[Scalar homogeneous dichotomy]
\label{thm:dichotomy}
Zero is fixed. For every nonzero seed \(z=x+y\theta\):
\begin{enumerate}
\item If \(|q|^a\ne|d|^b\) for all positive \(a,b\), then the full raw orbit tends to infinity in complex norm.
\item If \(|q|^a=|d|^b\) for some positive \(a,b\), then the full raw orbit is eventually periodic. Every coordinate satisfies
\[
|x_j|\le |q||d||x|,\qquad |y_j|\le |q||d||y|.
\]
\end{enumerate}
No restriction to nonreal seeds or to positive scalar coefficients is imposed.
\end{theorem}

\begin{proof}
First recall the positive content classification. A nonempty content return gives
\(|q|^A=|d|^E\), where \(A,E>0\) are its affine and division counts:
a word using only one branch cannot return a positive magnitude.
Under independence, a collision of two times would give such a return at the intervening positive state.
Thus the content orbit is injective. Each bounded subset of \(\N\) is finite, so an injective sequence of natural numbers tends to infinity.

If \(y\ne0\), Lemma~\ref{lem:coordinates} gives \(|y_j|\to\infty\), and
\[
|x_j+y_j\theta|\ge |\operatorname{Im}(x_j+y_j\theta)|
=\operatorname{Im}(\theta)|y_j|\longrightarrow\infty.
\]
If \(y=0\), this coordinate remains zero and \(x\ne0\).
The same lemma gives \(|x_j|\to\infty\), which is exactly complex-norm escape on the real axis.

In the dependent case, write \(|q|=t^u,\ |d|=t^v\), with \(t\ge2\), \(u,v>0\), and \(\gcd(u,v)=1\).
To obtain this form, take a dependence \(|q|^a=|d|^b\), let \(g=\gcd(a,b)\), and put \(u=b/g,\ v=a/g\).
For every prime \(\ell\), equality of exponents gives \(v_\ell(|q|)=u e_\ell\) and \(v_\ell(|d|)=v e_\ell\), since \(\gcd(u,v)=1\).
Then \(t=\prod_\ell\ell^{e_\ell}\) has the required properties and \(t\ge2\).
Write \(g_0=rt^k\) with \(t\nmid r\).
The residual \(r\) is fixed, and the exponent changes by
\(k\mapsto k-v\) when \(k\ge v\), and \(k\mapsto k+u\) otherwise.
Every exponent is less than \(\max(k+1,u+v)\):
above \(u+v-1\) the division branch strictly decreases, and below this bound both branches stay below it.
Therefore \(g_j\le t^{u+v}g_0=|q||d|g_0\).
Insert this bound into Lemma~\ref{lem:coordinates} and cancel the positive \(g_0\).
This gives the two asserted coordinate bounds.

The actual signed pair orbit lies in a finite integral rectangle.
Two states therefore coincide. Determinism makes every later state repeat at the same positive time difference, proving eventual periodicity of the raw pair.
The unique coordinate embedding recovers the actual ring states and their identical clocks.
\end{proof}

When both scalars are positive, the content exponent cycle is addition by \(u\) modulo \(u+v\), with primitive period \(u+v\).
For signed scalars the present theorem asserts eventual periodicity without a formula for its primitive period.
For example, \(q=4,d=8\) gives the nonreal cycle
\(\theta,4\theta,16\theta,2\theta,8\theta,\theta\).
The dependent signed classification concerns the actual finite rectangle, not just periodicity of the content factor.

\section{Formal correspondence}

The active file \texttt{QCD.\allowbreak Complex\allowbreak Lattice.\allowbreak 001/\allowbreak lean/\allowbreak Homogeneous\allowbreak Lattice.\allowbreak lean}
proves the content guard, its scalar recursion, the coordinate reconstruction identity, full-sequence height escape, and bounded dependent coordinates.
Its final declarations \texttt{Ring\allowbreak Independent\allowbreak Escape} and
\texttt{Ring\allowbreak Dependent\allowbreak Eventually\allowbreak Periodic} recover the actual quadratic ring and complex norm.
The positive content classification is the one in
\texttt{QCD.\allowbreak Homogeneous.\allowbreak 001} and \texttt{QCD.\allowbreak Homogeneous.\allowbreak 002}.
These are same-clock scalar dynamical statements; they do not identify the multiplication tables of different quadratic orders.

\end{document}
